\scanpage{15}
\deleted{\unsure{On|Le|Rien}} On a d'ailleurs
\[
\underset{\substack{\text{ens d'arêtes}\\\text{de }Q}}{A_Q}
 \simeq\underset{\substack{\text{produit}\\\text{des deux}\\\text{diagonales de }Q}}{\prod_{i\in\tau}D_i^*}
 \simeq V^*-\bigcup_{i\in\tau}D_i^*
 =\bigcup_{i\in\tau'}D_i^*=
\]
ce qui montre que \deleted{\unsure{$V^*$|$V$} \unsure{ou}} le carré défini par \deleted{\unsure{l'ens|l'ens.}}\inserted{la paire} \deleted{\unsure{des}} $(D_i)_{i\in\tau'}$ de \deleted{droites} de $V^*$ complémentaire de la paire $(D_i)_{i\in\tau}$ \unsure{d'identifie|s'identifie} au carré dual $DQ$ de $Q$. Si on désigne par $V(Q)$ le \inserted{plan} vectoriel sur $\mathbb F_3$ défini par $Q$, on trouve donc un isomorphisme
\[
V(DQ)\xrightarrow{\overset{\alpha_Q}{\sim}}V(Q)
\tag*{(11)}
\]
et de façon symétrique on a
\[
V\bigl(\underset{\substack{\rotatebox{90}{$\simeq$}\\DDQ}}{\deleted{\unsure{D}}Q}\bigr)
 \xrightarrow{\alpha_{DQ}}V(DQ)
\tag*{(\unsure{12|11})}
\]
Ceci posé, on a la relation
\[
\alpha_{\deleted{\unsure{D}}Q}\,\alpha_{DQ}=-\mathrm{id}_{V(Q)}
=V\bigl(\underset{\substack{\text{antipodisme}\\\text{de }Q}}{\underbrace{\underline{a}_{Q}}}\bigr)
\tag*{(\unsure{12|13})}
\]
et de $=$ bien entendu
\[
\alpha_{DQ}\,\alpha_Q=-\mathrm{id}_{V(DQ)}=V(\underline{a}_{DQ}).
\]
% archive sheet number: 14
